\section{Domain Task 10: Trải nghiệm khách hàng thông qua Khai phá văn bản}

Câu hỏi nghiệp vụ: Khách hàng thường đề cập đến những từ khóa nào nhất?
Những yếu tố cốt lõi nào tạo nên sự hài lòng (Điểm cao) và yếu tố nào dẫn
đến sự phàn nàn (Điểm thấp)?

Mục đích: Ứng dụng Khai phá văn bản (NLP) để lượng hóa cảm xúc từ bình
luận, giúp chủ nhà / nền tảng biết chính xác các điểm cần tối ưu.

\subsection{Biểu đồ 1: Đám mây từ vựng trải nghiệm (Word Cloud)}

\subsubsection*{A. Thiết kế Idiom}

\begin{table}[H]
\centering
\renewcommand{\arraystretch}{1.1}
\begin{tabularx}{\textwidth}{|l|l|X|}
\hline
\textbf{Idiom} & \multicolumn{2}{l|}{\textbf{Word Cloud (Đám mây từ vựng)}} \\ \hline
\textbf{What}
& \multicolumn{2}{X|}{Từ khóa (Keyword): Categorical (C)} \\ \cline{2-3}
& \multicolumn{2}{X|}{Tần suất xuất hiện (Count): Quantitative (Q)} \\ \hline
\textbf{How}
& \textbf{Encode} &
\textbf{Mark:} Text (Chữ viết) \newline
\textbf{Channel:} \newline
\hspace*{1em} Size (Kích thước chữ): Tỷ lệ thuận với tần suất đếm được \newline
\hspace*{1em} Color: Biểu diễn độ lớn của Tần suất (Q) (chữ càng đậm = xuất hiện càng nhiều) \\ \cline{2-3}
& \textbf{Reduce} &
Filter: Áp dụng tập stopwords mở rộng để loại bỏ các từ vô nghĩa. \\ \hline
\textbf{Why} & \multicolumn{2}{X|}{
\textbf{produce (derive) $\rightarrow$ explore $\rightarrow$ summarize} \newline
Dẫn xuất (derive): bóc tách văn bản loại bỏ stopwords để lấy ra các unigram/bigram (1--2 từ) có ý nghĩa. Explore: Khám phá nhanh và tóm tắt (summarize) các chủ đề / ngữ cảnh lưu trú được thảo luận nhiều nhất.
} \\ \hline
\textbf{Scale} & \multicolumn{2}{X|}{
Main key: Top các từ khóa xuất hiện nhiều nhất.
} \\ \hline
\end{tabularx}
\end{table}

\begin{figure}[H]
    \centering
    \safeincludegraphics[width=0.9\linewidth]{charts/domain_task10_chart1.png}
    \caption{Hình 10.1. Biểu đồ đám mây từ vựng trải nghiệm của khách hàng}
    \label{fig:domain-task10-chart1}
\end{figure}

\subsubsection*{B. Đánh giá biểu đồ}

\textbf{1. Nguyên lý biểu đạt (Expressiveness):}

\begin{itemize}
    \item Thuộc tính: Keyword (C), Count (Q).
    \item Channel: Sử dụng Mark là Text để trực tiếp hiển thị thuộc tính Keyword (C). Sử dụng Channel Kích thước (Size) và Channel Độ đậm nhạt (Color Luminance) để cùng biểu diễn độ lớn định lượng (Q).
    \item Đánh giá mức độ quan trọng:
    \begin{itemize}
        \item Ưu tiên 1 -- Kích thước (Size/Area): Thuộc tính định lượng (Tần suất xuất hiện -- Q) được gán vào Size. Theo Mackinlay, Area/Size là kênh yếu đối với dữ liệu (Q), nhưng trong đặc thù của Word Cloud, nó là yếu tố duy nhất tạo ra hiệu ứng nổi bật cho các từ khóa chính.
        \item Ưu tiên 2 -- Độ đậm nhạt màu (Color Luminance): Vì kênh Size bị ảnh hưởng bởi độ dài của text (chữ dài trông to hơn chữ ngắn dù cùng tần suất), kênh Color Luminance được bổ sung để biểu diễn Tần suất (Q). Theo Mackinlay, Luminance tuy xếp hạng thấp cho (Q) nhưng ở đây nó hỗ trợ đắc lực cho kênh Size: Từ nào vừa to vừa đậm (như stay, clean, great) chắc chắn là từ phổ biến nhất.
    \end{itemize}
\end{itemize}

\textbf{2. Nguyên lý hiệu quả (Effectiveness):}

\begin{itemize}
    \item Độ chính xác (Accuracy): Kênh Area/Size có độ chính xác rất thấp. Mắt người có xu hướng ước lượng sai lệch diện tích của chữ cái, bị nhiễu bởi độ dài của từ (ví dụ chữ ``recommendation'' tự nhiên trông to hơn chữ ``nice'' dù cùng count). Log\_error rất lớn, không phù hợp để so sánh định lượng khắt khe.
    \item Khả năng phân biệt (Discriminability): Biểu đồ có hiện tượng lộn xộn (clutter). Chỉ nổi bật được vài từ vựng lớn nhất ở trung tâm, các từ nhỏ xung quanh rất khó đọc và khó phân biệt cấp độ.
    \item Khả năng tách biệt (Separability): Kênh vị trí (Pos) và kênh màu sắc (Color) tách biệt hoàn toàn. Việc các điểm nằm ở vị trí cao/thấp không làm ảnh hưởng đến khả năng nhận diện màu sắc của quận đó.
\end{itemize}

\subsubsection*{C. Phân tích biểu đồ (Insight)}

\begin{itemize}
    \item Nhìn tổng quan, trải nghiệm của khách hàng xoay quanh 3 trụ cột chính: Vị trí (location, subway, close), Không gian (clean, comfortable, apartment) và Chủ nhà (host, responsive, helpful).
\end{itemize}
